% GENERATED FILE. Edit canonical lessons, then run npm run generate:books.
% canonical-source-hash: fnv1a64:a4594050c16d0386
% canonical-lessons: GE-C370-klappen, GE-C370-misslingen, GE-C370-eintreten, GE-C370-zutreffen, GE-C370-herausstellen

\chapter{Ob es klappt}
\label{ch:ge-ob-es-klappt}
\hlchaptermodality{370}

\begin{chapteropening}
\textbf{By the end of this chapter:} \emph{I can say whether something worked out or went wrong, whether a case arose and whether a claim is true.}

Complete the last lesson of chapter 370: I can say whether something worked out or went wrong, whether a case arose and whether a claim is true.
\end{chapteropening}

\section[klappen]{klappen — to work out (informal); to fold}

\label{lesson:GE-C370-klappen}



\begin{quote}
Before the new one: say the German for the guarantee, then the German for to guarantee.
\end{quote}

\subsection*{You'll want to know: klappen}
\textbf{klappen} — \textquotedblleft{}to work out (informal); to fold\textquotedblright{}.

\begin{quote}
Mit der neuen Wohnung hat es endlich geklappt.
\end{quote}

\begin{quote}
\emph{It finally worked out with the new flat.}
\end{quote}

Informal and very common: \emph{Das klappt nicht} — that does not work. Literally to fold or flap: \emph{den Sitz nach vorne klappen}.

\begin{cousinweb}
\emph{klappen} imitates a sharp sound; EtymDB lists English \emph{clap} and Dutch \emph{klappen} as cousins. It is usually said to have come into the standard language from Low German. The sense \textquotedblleft{}to work out\textquotedblright{} is usually explained as things snapping neatly into place.
\end{cousinweb}

\begin{grammarlens}[title={What you've built}]
One more word for telling, explaining and describing at B1.
\end{grammarlens}

\subsection*{Your turn}
\begin{itemize}
\raggedright
  \item \emph{Say it:} \emph{klappen}
  \item \emph{Say it:} \emph{klappen}, once more
  \item \emph{Say it:} a sentence of your own with \emph{klappen}
  \item \emph{Recall:} say the German for the guarantee, then the German for to guarantee, then say \emph{klappen} again
\end{itemize}

\begin{tcolorbox}[breakable,colback=teal!4,colframe=teal!35!black,title={Before you move on}]
What does klappen mean? (\textquotedblleft{}To work out (informal); to fold\textquotedblright{}.) Say it once more.
\end{tcolorbox}
\section[misslingen]{misslingen — to go wrong, to fail (of things)}

\label{lesson:GE-C370-misslingen}



\begin{quote}
Before the new one: say the German for to guarantee, then the German for to work out.
\end{quote}

\subsection*{You'll want to know: misslingen}
\textbf{misslingen} — \textquotedblleft{}to go wrong, to fail (of things)\textquotedblright{}.

\begin{quote}
Der erste Kuchen ist mir misslungen.
\end{quote}

\begin{quote}
\emph{My first cake went wrong.}
\end{quote}

Perfekt with \emph{sein}: \emph{es ist misslungen}. The thing is the subject and the person goes in the dative: \emph{Es misslingt mir.}

\begin{cousinweb}
\emph{misslingen} is built with \emph{miss-} \textquotedblleft{}wrongly\textquotedblright{} on the old verb behind \emph{gelingen} \textquotedblleft{}to succeed\textquotedblright{}, Old High German \emph{gilingan}, Proto-Germanic \emph{lingwaną} (EtymDB).
\end{cousinweb}

\begin{grammarlens}[title={What you've built}]
One more word for telling, explaining and describing at B1.
\end{grammarlens}

\subsection*{Your turn}
\begin{itemize}
\raggedright
  \item \emph{Say it:} \emph{misslingen}
  \item \emph{Say it:} \emph{misslingen}, once more
  \item \emph{Say it:} a sentence of your own with \emph{misslingen}
  \item \emph{Recall:} say the German for to guarantee, then the German for to work out, then say \emph{misslingen} again
\end{itemize}

\begin{tcolorbox}[breakable,colback=teal!4,colframe=teal!35!black,title={Before you move on}]
What does misslingen mean? (\textquotedblleft{}To go wrong, to fail (of things)\textquotedblright{}.) Say it once more.
\end{tcolorbox}
\section[eintreten]{eintreten — to happen, to occur; to enter}

\label{lesson:GE-C370-eintreten}



\begin{quote}
Before the new one: say the German for to work out, then the German for to go wrong.
\end{quote}

\subsection*{You'll want to know: eintreten}
\textbf{eintreten} — \textquotedblleft{}to happen, to occur; to enter\textquotedblright{}.

\begin{quote}
Der schlimmste Fall ist zum Glück nicht eingetreten.
\end{quote}

\begin{quote}
\emph{Luckily the worst case did not happen.}
\end{quote}

Separable, Perfekt with \emph{sein}: \emph{es tritt ein}, \emph{es ist eingetreten}. Also to enter: \emph{in einen Verein eintreten} — to join a club.

\begin{cousinweb}
\emph{eintreten} is built from \emph{ein} \textquotedblleft{}in\textquotedblright{} and \emph{treten} \textquotedblleft{}to step\textquotedblright{}, inherited from Old High German \emph{tretan} and Proto-Germanic \emph{tredaną} (EtymDB): an event steps in.
\end{cousinweb}

\begin{grammarlens}[title={What you've built}]
One more word for telling, explaining and describing at B1.
\end{grammarlens}

\subsection*{Your turn}
\begin{itemize}
\raggedright
  \item \emph{Say it:} \emph{eintreten}
  \item \emph{Say it:} \emph{eintreten}, once more
  \item \emph{Say it:} a sentence of your own with \emph{eintreten}
  \item \emph{Recall:} say the German for to work out, then the German for to go wrong, then say \emph{eintreten} again
\end{itemize}

\begin{tcolorbox}[breakable,colback=teal!4,colframe=teal!35!black,title={Before you move on}]
What does eintreten mean? (\textquotedblleft{}To happen, to occur; to enter\textquotedblright{}.) Say it once more.
\end{tcolorbox}
\section[zutreffen]{zutreffen — to be true, to apply}

\label{lesson:GE-C370-zutreffen}



\begin{quote}
Before the new one: say the German for to go wrong, then the German for to happen.
\end{quote}

\subsection*{You'll want to know: zutreffen}
\textbf{zutreffen} — \textquotedblleft{}to be true, to apply\textquotedblright{}.

\begin{quote}
Das trifft nicht auf alle Schüler zu.
\end{quote}

\begin{quote}
\emph{That does not apply to all pupils.}
\end{quote}

Separable and strong: \emph{das trifft zu}, \emph{das hat zugetroffen}. \emph{Zutreffend} — correct.

\begin{cousinweb}
\emph{zutreffen} is built from \emph{zu} and \emph{treffen} \textquotedblleft{}to hit\textquotedblright{}, inherited from Old High German \emph{treffan} and Proto-Germanic \emph{drepaną} (EtymDB): a statement that hits the mark.
\end{cousinweb}

\begin{grammarlens}[title={What you've built}]
One more word for telling, explaining and describing at B1.
\end{grammarlens}

\subsection*{Your turn}
\begin{itemize}
\raggedright
  \item \emph{Say it:} \emph{zutreffen}
  \item \emph{Say it:} \emph{zutreffen}, once more
  \item \emph{Say it:} a sentence of your own with \emph{zutreffen}
  \item \emph{Recall:} say the German for to go wrong, then the German for to happen, then say \emph{zutreffen} again
\end{itemize}

\begin{tcolorbox}[breakable,colback=teal!4,colframe=teal!35!black,title={Before you move on}]
What does zutreffen mean? (\textquotedblleft{}To be true, to apply\textquotedblright{}.) Say it once more.
\end{tcolorbox}
\section[herausstellen]{herausstellen — to turn out (reflexive, sich herausstellen)}

\label{lesson:GE-C370-herausstellen}



\begin{quote}
Before the new one: say the German for to happen, then the German for to be true.
\end{quote}

\subsection*{You'll want to know: herausstellen}
\textbf{herausstellen} — \textquotedblleft{}to turn out (reflexive, sich herausstellen)\textquotedblright{}.

\begin{quote}
Später hat sich herausgestellt, dass er recht hatte.
\end{quote}

\begin{quote}
\emph{Later it turned out that he was right.}
\end{quote}

Reflexive and separable: \emph{es stellt sich heraus, dass …}, \emph{es hat sich herausgestellt}.

\begin{cousinweb}
\emph{herausstellen} is built from \emph{heraus} \textquotedblleft{}out\textquotedblright{} and \emph{stellen} \textquotedblleft{}to put\textquotedblright{}, Old High German \emph{stellen}, Proto-Germanic \emph{stallijaną} (EtymDB): the truth is put out where all can see it.
\end{cousinweb}

\begin{grammarlens}[title={What you've built}]
One more word for telling, explaining and describing at B1.
\end{grammarlens}

\subsection*{Your turn}
\begin{itemize}
\raggedright
  \item \emph{Say it:} \emph{herausstellen}
  \item \emph{Say it:} \emph{herausstellen}, once more
  \item \emph{Say it:} a sentence of your own with \emph{herausstellen}
  \item \emph{Recall:} say the German for to happen, then the German for to be true, then say \emph{herausstellen} again
\end{itemize}

\begin{tcolorbox}[breakable,colback=teal!4,colframe=teal!35!black,title={Before you move on}]
What does herausstellen mean? (\textquotedblleft{}To turn out (reflexive, sich herausstellen)\textquotedblright{}.) Say it once more.
\end{tcolorbox}
